\documentclass[aspectratio=169]{beamer}

% ===== THEME CHO MÁY CHIẾU =====
\usetheme{CambridgeUS}
\usecolortheme{default}
\setbeamertemplate{navigation symbols}{}

% ===== FONT & TIẾNG VIỆT (XeLaTeX) =====
\usepackage{fontspec}
\usepackage{polyglossia}
\setdefaultlanguage{vietnamese}
\setmainfont{Latin Modern Roman}

% ===== CODE =====
\usepackage{listings}
\lstset{
  basicstyle=\ttfamily\Large,
  columns=fullflexible,
  frame=single
}

% ===== TITLE =====
\title{\Huge STACK \\[6pt] \Large (Ngăn xếp)}
\subtitle{\Large Tư duy giải quyết vấn đề}
\author{\large Giáo viên lập trình}
\date{}

\begin{document}

%------------------------------------------------
\begin{frame}
  \titlepage
\end{frame}

%------------------------------------------------
\begin{frame}{Mục tiêu bài học}
\Large
\begin{itemize}
    \item Hiểu \textbf{ý tưởng cốt lõi} của Stack
    \item Biết \textbf{khi nào nên dùng}
    \item Tự giải được bài toán tương tự
\end{itemize}
\end{frame}

%------------------------------------------------
\begin{frame}{Một tình huống quen thuộc}
\Large
\begin{itemize}
    \item Một \textbf{chồng sách} trên bàn
    \item Lấy sách: luôn lấy \textbf{cuốn trên cùng}
    \item Cuốn đặt \textbf{sau} → lấy ra \textbf{trước}
\end{itemize}
\end{frame}

%------------------------------------------------
\begin{frame}{Quan sát quan trọng}
\Large
\begin{block}{}
\centering
Chỉ thao tác được ở \textbf{một đầu}
\end{block}

\vspace{10pt}

\begin{block}{}
\centering
Thứ vào sau lại được xử lý trước
\end{block}
\end{frame}

%------------------------------------------------
\begin{frame}{Ý tưởng cốt lõi của Stack}
\Huge
\begin{alertblock}{}
\centering
VÀO SAU -- RA TRƯỚC
\end{alertblock}
\end{frame}

%------------------------------------------------
\begin{frame}{Ví dụ: Chuỗi ngoặc}
\Huge
\centering
\texttt{( ( ) ( ) )}
\end{frame}

%------------------------------------------------
\begin{frame}{Cách suy nghĩ}
\Large
\begin{itemize}
    \item Gặp ``('' → \textbf{chưa xử lý}
    \item Gặp ``)'' → ghép với ``('' \textbf{gần nhất}
\end{itemize}

\vspace{12pt}
\pause
\centering
\textbf{Làm sao nhớ được dấu ``('' gần nhất?}
\end{frame}

%------------------------------------------------
\begin{frame}{Diễn giải trực quan}
\Large
\begin{center}
\begin{tabular}{c|c}
Ký tự & Đang giữ \\
\hline
( & ( \\
( & ( ( \\
) & ( \\
( & ( ( \\
) & ( \\
) & rỗng \\
\end{tabular}
\end{center}

\vspace{10pt}
\centering
\textbf{→ Chuỗi hợp lệ}
\end{frame}

%------------------------------------------------
\begin{frame}{Chuyển sang lập trình}
\Large
\begin{itemize}
    \item Giữ lại → \textbf{push}
    \item Lấy ra → \textbf{pop}
    \item Kiểm tra rỗng → \textbf{empty}
\end{itemize}

\vspace{10pt}
\centering
Stack chỉ thao tác ở \textbf{đỉnh}
\end{frame}

%------------------------------------------------
\begin{frame}{Code C++ minh họa}
\begin{lstlisting}[language=C++]
stack<char> st;

for (char c : s) {
    if (c == '(')
        st.push(c);
    else {
        if (st.empty()) {
            cout << "Sai";
            return 0;
        }
        st.pop();
    }
}

cout << (st.empty() ? "Dung" : "Sai");
\end{lstlisting}
\end{frame}

%------------------------------------------------
\begin{frame}{Bài tập luyện tập}
\Large
\textbf{Bài 1:}
\begin{itemize}
    \item Chuỗi chỉ có ( và )
\end{itemize}

\vspace{10pt}
\textbf{Bài 2:}
\begin{itemize}
    \item Có thêm [ và ]
\end{itemize}

\vspace{12pt}
\centering
\textbf{Gợi ý: dấu mở gần nhất}
\end{frame}

%------------------------------------------------
\begin{frame}{Tổng kết}
\Large
\begin{itemize}
    \item Dùng Stack khi cần \textbf{ngược thứ tự}
    \item Phù hợp bài toán ghép cặp
    \item Tư duy đúng → code rất ngắn
\end{itemize}
\end{frame}

\end{document}
